\documentclass[aspectratio=169]{beamer}

\usetheme{Madrid}
\usepackage{booktabs}
\usepackage{xcolor}
\setbeamertemplate{navigation symbols}{}

\title[Self-explaining LULC workbench]{From classifier to product: a self-explaining LULC workbench, a richer model zoo, and portable outputs}
\author[Salil Gujar]{Salil Gujar \\ \footnotesize M.Tech CSE \quad Entry No.\ 2025MCS2106}
\date{July 2026}

\begin{document}

\frame{\titlepage}

% ------------------------------------------------------------------
\begin{frame}{Where we are, and what week 8 tackled}
\begin{columns}[T]
\begin{column}{0.5\textwidth}
\emph{Done so far}
\begin{itemize}\small
  \item A 4-class base map over India at 10\,m with a living hierarchy: split or add a class, give a
  few example polygons, and retrain that node on the fly.
  \item A git-backed zoo of Model and Dataset cards, with save, reload, merge, base and year
  selection, publishing, and site testing on acacia, tea, and mining.
\end{itemize}
\end{column}
\begin{column}{0.5\textwidth}
\emph{What was advised (26 points)}
\begin{itemize}\small
  \item Make it explainable and controllable: show only the controls that fit the flow, and guide the
  user through the pipeline.
  \item Make the zoo honest: standard classes, dataset-to-model links, selective publishing.
  \item Add new abilities: Tessera training, an auto model bake-off, a GeoTIFF export, a resumable
  project.
  \item Answer the operational questions on schemas, deployment, and running Tessera with joblib.
\end{itemize}
\end{column}
\end{columns}
\end{frame}

% ------------------------------------------------------------------
\begin{frame}{A workbench that explains itself}
\begin{itemize}
  \item Two panels, one job each. The left panel is navigation (area, hierarchy, save, zoo); a new
  right panel is contextual: click a class and it shows exactly what you can do with it, in order:
  mark data, split, add, merge, retrain.
  \item Only the relevant controls. The balancing and multi-year knobs appear only when you are
  training your own split, not for a bare leaf or the base map.
  \item Two ways to see the scheme. ``By hierarchy'' is the class tree; ``by operations'' is the
  ordered sequence that built it, and clicking a step acts on that class. The same right panel serves
  both views.
  \item Draw the area on the map, with a rectangle tool or a click-to-centre size slider; the
  classification is clipped to the box, and an eye toggle hides or shows it.
  \item Moving to a new area offers a clean reset, confirmed first, so a scheme built for one site does
  not leak into the next.
\end{itemize}
\end{frame}

% ------------------------------------------------------------------
\begin{frame}{Model zoo improvements}
\begin{itemize}
  \item Every trained model shows up as a card. We fixed the case where re-splitting a node dropped
  its old model, so nothing trained goes missing, and any card can be deleted.
  \item Standard classes on the card. If the uploader maps a class to a standard scheme (ESA
  WorldCover or USDA), the small card shows the standard name, for example ``Tree cover''; the detail
  pane still shows what the uploader calls it, acacia mapped to Tree cover.
  \item Provenance both ways. A dataset card now lists which models use it, alongside the existing
  model-to-dataset links, so one dataset feeding several models is visible from either side.
  \item Selective publishing. Tick the cards you want and publish those, instead of only ``publish
  all''.
\end{itemize}
\end{frame}

% ------------------------------------------------------------------
\begin{frame}{Auto model bake-off}
``Auto'' trains several models on the same held-out-polygon split and keeps the most accurate. We
restrict the field to linear models, because only a linear model replays as Earth-Engine band math, so
the winner still renders as the crisp 10\,m tile map.
\vspace{4mm}
\centering
\begin{tabular}{lc}
\toprule
Candidate (acacia / non-acacia) & Held-out accuracy \\
\midrule
LinearSVC & 0.703 \\
Logistic Regression & 0.692 \\
Ridge & 0.731 \\
\bottomrule
\end{tabular}
\end{frame}

% ------------------------------------------------------------------
\begin{frame}{Tessera as a training embedding}
\begin{itemize}\small
  \item A user can now train a split on Tessera (128-d) instead of Alpha Earth (64-d), for example
  acacia against non-acacia. It is scored, carded, and kept like any other model.
  \item On a Delhi run, acacia against non-acacia on Tessera reached about 0.73 held-out, comparable to
  Alpha Earth on the same split.
  \item Our four requested ROIs were approved for 2017 to 2025 (v1.1, cambridge variant). Because only
  these four boxes were approved, we have the multi-year Tessera data only inside those boxes, so
  Tessera training is offered only there.
  \item A Tessera split is a per-point download rather than an Earth-Engine image, so it does not ride
  the crisp tile map. It is evaluated and carded, and the card is marked accordingly.
\end{itemize}
\end{frame}

% ------------------------------------------------------------------
\begin{frame}{Portable outputs: GeoTIFF and a resumable project}
\begin{itemize}
  \item Download the result as a GeoTIFF. The classified area exports as a single-band integer-class
  raster at 10\,m, with a code-to-class legend, built server-side and fetched by the browser, ready to
  open in QGIS.
  \item Save the whole project as one JSON: the class scheme, the sequence of steps that built it, and
  the area, year, and base classes. Datasets and models ride as links, never copied. Reloading resumes
  the exact view.
\end{itemize}
\end{frame}

% ------------------------------------------------------------------
\begin{frame}{New ground truth: seasonal water}
We downloaded a set of water and non-water polygons for farm ponds and water bodies, marked across
about ten dates each between 2019 and 2025, and folded them into the zoo as a dataset card. This is a
seasonal signal: water present or absent by date, keyed by a per-water-body id.

\vspace{4mm}
\begin{center}
\begin{tabular}{lr}
\toprule
What we have & Count \\
\midrule
Total polygons (WGS84) & 876 \\
\quad water & 720 \\
\quad non-water & 156 \\
Spatial spread (0.25$^\circ$ grid) & 0.61 \\
\bottomrule
\end{tabular}
\end{center}
\end{frame}

% ------------------------------------------------------------------
\begin{frame}{Deploying this on a server}
\begin{itemize}
  \item Runtime: a Python virtual environment from the requirements file, served with uvicorn. State
  is plain files (the hierarchy, the cards, the model binaries), so there is no database to run.
  \item Earth Engine: the app initialises with a project id and the runtime user's Earth Engine
  credentials, so on a server we authenticate once as that user. A service-account key is the natural
  choice for a fully headless box.
  \item Publishing to GitHub: the model zoo is its own git repository pushed to a configured remote.
  We leave authentication to the system git, so a server needs either a token through a credential
  helper or an SSH deploy key.
  \item Configuration is a small set of environment variables: the Earth-Engine project and user, and
  the zoo remote. Committed model binaries let a fresh clone classify straight away.
\end{itemize}
\end{frame}

% ------------------------------------------------------------------
\begin{frame}
\centering
\vspace{1.5cm}
{\Large Thank you}
\vspace{0.5cm}


\end{frame}

\end{document}
